\documentclass[10pt,a4paper]{amsart}
\usepackage[margin=19mm]{geometry}
\usepackage{amsmath,amssymb,mathtools}
\usepackage[scr=rsfs,cal=euler]{mathalfa}
\usepackage{xfrac}
\usepackage[unicode,hidelinks,bookmarks=false]{hyperref}
\newcommand{\ie}{i.e.\ }
\newcommand{\cf}{cf.\ }
\newcommand{\resp}{respectively\ }
\newcommand{\Subsection}{\subsection}
\providecommand{\nobreakdash}{\nobreak\hbox{-}}
\setlength{\emergencystretch}{2em}
\multlinegap0pt
\usepackage{amssymb}
\usepackage{xcolor}
\newtheorem{theorem}{Theorem}
\newtheorem{corollary}{Corollary}
\newcommand{\CS}{\mathcal{S}}
\newcommand{\Vol}{\operatorname{vol}_g}
\title{Differentiability and other properties of the cosmological volume function}
\author{Leonardo Garc\'ia-Heveling}
\date{Source: arXiv:2512.11626v2 (CC BY 4.0)}
\begin{document}
\maketitle

\noindent\emph{Source and licence.} Differentiability and other properties of the cosmological volume function. Version arXiv:2512.11626v2 (CC BY 4.0), distributed by arXiv under the Creative Commons Attribution 4.0 International licence, http://creativecommons.org/licenses/by/4.0/, as recorded in the arXiv licence metadata for this exact version and shown on the abstract page https://arxiv.org/abs/2512.11626v2. Full licence notice: \texttt{licenses/arxiv-2512.11626-CC-BY-4.0.txt}.\par\smallskip

\noindent\textit{Editorial scope.} Counted unit: the article's corollary cor:2 (source0.tex lines 287--289). The article's complete original proof of it is delivered with it (source0.tex lines 291--301, one \texttt{proof} environment of 11 lines). No mathematics is written, repaired or re-derived: the statement and the proof are the article's own lines, in the article's own notation, and no retained line is substituted or re-flowed.  Retained uncounted as the source-local context the proof needs: lines 64--66.  Nothing was omitted from any retained span, so the package carries no omission record.  No retained line contains a citation, so the wrapper carries no bibliography and no bibliography entry.  No retained line contains a figure environment, an image-inclusion command or drawing code, so no image is delivered and the package declares no asset dependency.  Editorial formatting only, in addition: an AMS \texttt{amsart} preamble that restates the article's own declarations and macros used by the retained lines, namely the environments \texttt{theorem}, \texttt{corollary} on separate counters, together with the standard packages those lines need.  The retained environments print in this document as Theorem 1, Corollary 1 in order of appearance, whereas the article prints the same items with the numbering of its own full document.  The complete native source of the article as distributed in the arXiv e-print for this exact version is bundled unchanged as \texttt{sources/190887/source0.tex}.
\begin{theorem} \label{thm:main}
    Let $(M,g)$ be an orientable spacetime and $\CS$ a future Cauchy surface. Then the cosmological volume function $\tau_V$ of the spacetime $(I^+(\CS),g)$ is regular and is a $C^1$ temporal function.
\end{theorem}

\begin{corollary} \label{cor:2}
    Let $(M,g)$ be a spacetime with regular cosmological volume function $\tau_V$. If $p \in M$ is a point such that $I^-(p)$ contains a future Cauchy surface $\CS$, then $\tau_V$ is $C^1$ with timelike gradient in a neighborhood of $p$.
\end{corollary}

\begin{proof}
    Since $I^-(p) \cap \CS = \CS$ is compact, and chronological pasts are open, it is easy to infer the existence of a neighborhood $U$ of $p$ such that $\CS \subset I^-(q)$ for every $q \in U$. Then
    \begin{equation*}
        I^-(q) = \left(I^-(q) \cap I^+(\CS)\right) \cup \CS \cup I^-(\CS),
    \end{equation*}
    which can be seen as follows: For each $x \in I^-(q)$, let $\gamma$ be the timelike curve $\gamma$ joining $x$ and $q$. Either $\gamma$ intersects $\CS$, in which case $x \in \CS \cup I^-(\CS)$, or it doesn't, in which case $x \in I^+(\CS)$, because $\gamma$ can be extended to the past until it intersects $\CS$. Then
    \begin{equation*}
        \tau_V(q) = \Vol\left(I^-(q) \cap I^+(\CS)\right) + \Vol\left(I^-(\CS)\right),
    \end{equation*}
    where the first summand is $C^1$ by Theorem~\ref{thm:main}, and the second summand is constant.
\end{proof}

\end{document}
